% ===== PRÉAMBULE CANONIQUE — à coller VERBATIM en tête de CHAQUE .tex =====
\documentclass[11pt,a4paper]{article}
\usepackage[utf8]{inputenc}\usepackage[T1]{fontenc}\usepackage{lmodern}
\usepackage[french]{babel}
\usepackage{amsmath,amssymb,mathtools}\usepackage{icomma}
\usepackage[version=4]{mhchem}          % \ce{...} : équations & formules
\usepackage{siunitx}\sisetup{locale=FR} % \qty{}{}, \si{}, \ang{}
\usepackage{chemfig}                    % \chemfig{...} : structures développées
\usepackage{xcolor}\usepackage{colortbl}
\usepackage{tikz}\usepackage{pgfplots}\pgfplotsset{compat=1.18}
\usepackage[most]{tcolorbox}\usepackage{enumitem}\usepackage{array,booktabs}
\usepackage[a4paper,margin=2.2cm]{geometry}
\usetikzlibrary{arrows.meta,calc,angles,quotes,positioning,patterns,backgrounds,shapes.geometric}
\definecolor{primary}{RGB}{222,49,99}\definecolor{primarydark}{RGB}{150,22,55}
\definecolor{defcol}{RGB}{0,114,178}\definecolor{methcol}{RGB}{52,92,148}
\definecolor{excol}{RGB}{0,135,90}\definecolor{attncol}{RGB}{200,40,55}
\tcbset{boxbase/.style={enhanced,breakable,boxrule=0.9pt,arc=2.5pt,
  left=3mm,right=3mm,top=2.3mm,bottom=2.3mm,fonttitle=\bfseries,coltitle=white}}
% --- boîtes TUTORIEL (noms EXACTS, ne pas renommer) ---
\newtcolorbox{cle}[1][]{boxbase,colback=defcol!6,colframe=defcol,
  title={\textsc{À retenir}\ifx\relax#1\relax\relax\else\textmd{ : \itshape #1}\fi}}
\newtcolorbox{methode}[1][]{boxbase,colback=methcol!7,colframe=methcol,sharp corners,
  title={\textsc{Méthode}\ifx\relax#1\relax\relax\else\textmd{ --- \itshape #1}\fi}}
\newtcolorbox{exemplebox}[1][]{boxbase,colback=excol!5,colframe=excol,
  title={\textsc{Exemple}\ifx\relax#1\relax\relax\else\textmd{ : \itshape #1}\fi}}
\newtcolorbox{piege}[1][]{boxbase,colback=attncol!6,colframe=attncol,
  title={\textsc{Piège}\ifx\relax#1\relax\relax\else\textmd{ : \itshape #1}\fi}}
\newtcolorbox{remarque}[1][]{boxbase,colback=black!4,colframe=black!45,
  title={\textsc{Remarque}\ifx\relax#1\relax\relax\else\textmd{ : \itshape #1}\fi}}
% --- boîtes EXERCICE / CORRIGÉ (noms EXACTS, auto-comptées) ---
\newtcolorbox[auto counter]{exo}[1][]{boxbase,colback=primary!4,colframe=primary,
  title={\textsc{Exercice}~\thetcbcounter\ifx\relax#1\relax\relax\else\textmd{ --- \itshape #1}\fi}}
\newtcolorbox[auto counter]{corrige}[1][]{boxbase,colback=excol!4,colframe=excol,
  title={\textsc{Corrigé}~\thetcbcounter\ifx\relax#1\relax\relax\else\textmd{ --- \itshape #1}\fi}}
\newcommand{\R}{\mathbb{R}}\newcommand{\N}{\mathbb{N}}\newcommand{\Z}{\mathbb{Z}}\newcommand{\ds}{\displaystyle}
% (un \usepackage{fancyhdr}+en-tête est possible ; il est ignoré côté web)
% =========================================================================
\begin{document}

\begin{exo}[identifier un ion par sa composition]
\textit{Données :} $Z(\mathrm{Cr})=24$, $Z(\mathrm{Fe})=26$, $Z(\mathrm{Ni})=28$, $Z(\mathrm{Zn})=30$.

\smallskip
Parmi les quatre cations ci-dessous, un seul possède \textbf{exactement $30$ neutrons et $24$ électrons}.
Détermine la composition $\big(n(p^+),n(n^0),n(e^-)\big)$ de chacun, puis désigne le bon.
\begin{enumerate}[label=\alph*)]
  \item $\ce{^{52}_{24}Cr^{2+}}$
  \item $\ce{^{56}_{26}Fe^{2+}}$
  \item $\ce{^{58}_{28}Ni^{2+}}$
  \item $\ce{^{64}_{30}Zn^{2+}}$
\end{enumerate}
\end{exo}

\begin{exo}[comparer neutrons et électrons]
\textit{Données :} $Z(\mathrm{N})=7$, $Z(\mathrm{Mg})=12$, $Z(\mathrm{Cl})=17$, $Z(\mathrm{K})=19$.

\smallskip
Pour chaque espèce, donne le nombre de neutrons et le nombre d'électrons, puis réponds : pour
\emph{laquelle} le nombre d'électrons est-il \textbf{égal} au nombre de neutrons ?
\begin{enumerate}[label=\alph*)]
  \item $\ce{^{14}_{7}N}$
  \item $\ce{^{24}_{12}Mg^{2+}}$
  \item $\ce{^{37}_{17}Cl^{-}}$
  \item $\ce{^{39}_{19}K^{+}}$
\end{enumerate}
\end{exo}

\begin{exo}[charge d'un grand nombre de particules]
On donne la charge élémentaire $e = \SI{1,6e-19}{\coulomb}$. Travaille avec les puissances de dix.
\begin{enumerate}[label=\alph*)]
  \item Quelle est la charge d'un ensemble de $2 \times 10^{13}$ protons ?
  \item Une gouttelette a capté $5 \times 10^{10}$ électrons en excès. Quelle est sa charge ?
  \item Combien de protons faut-il rassembler pour porter une charge totale de
        $+\SI{1,6e-7}{\coulomb}$ ?
\end{enumerate}
\end{exo}

\begin{exo}[mise en situation : le chlorure de calcium]
Le chlorure de calcium $\ce{CaCl2}$ est un solide ionique formé d'ions calcium et chlorure. On
considère les isotopes $\ce{^{40}_{20}Ca}$ et $\ce{^{35}_{17}Cl}$.
\begin{enumerate}[label=\alph*)]
  \item Donne la composition $\big(n(p^+),n(n^0),n(e^-)\big)$ de l'ion calcium $\ce{Ca^{2+}}$.
  \item Donne la composition de l'ion chlorure $\ce{Cl^{-}}$.
  \item Dans un motif $\ce{CaCl2}$ (un $\ce{Ca^{2+}}$ et deux $\ce{Cl^{-}}$), compte le nombre
        \emph{total} de protons et le nombre \emph{total} d'électrons. Vérifie que le motif est
        électriquement neutre.
\end{enumerate}
\end{exo}

\begin{exo}[vrai ou faux — synthèse]
\textit{Données :} $Z(\mathrm{Ar})=18$, $Z(\mathrm{K})=19$, $Z(\mathrm{Ca})=20$.

\smallskip
Pour chaque proposition, indique si elle est \textbf{vraie} ou \textbf{fausse}, en justifiant d'une
phrase.
\begin{enumerate}[label=\alph*)]
  \item Le nombre de masse $A$ est la somme du nombre de protons et du nombre de neutrons du noyau.
  \item En passant de l'atome $\ce{^{40}_{20}Ca}$ à l'ion $\ce{Ca^{2+}}$, le nombre de neutrons diminue.
  \item L'ion $\ce{^{40}_{20}Ca^{2+}}$ possède le même nombre d'électrons que l'atome $\ce{^{40}_{18}Ar}$.
  \item Dans le noyau de l'ion $\ce{^{39}_{19}K^{+}}$, il y a $18$ électrons.
\end{enumerate}
\end{exo}

\end{document}
